\section{An exact renormalization step with almost dimension-two \texorpdfstring{\mbox{contraction}}{contraction}}
\label{sec:rg}

The linear operators of Section~\ref{sec:free} identify the quadratic
part of one block integration.  We now control the nonlinear remainder,
including configurations with large spin differences.  Two features of
$S^2$ matter.  Its tangent planes do not have a globally preferred frame;
we therefore perform Gaussian completion in the fixed ambient space
$\R^3$.  In addition, the full stabilizer $O(2)$ contains tangent
inversion.  Together with spatial inversion, this removes the cubic
terms that would otherwise limit the contraction of the normalized
remainder to order $L^{-1}$.

Throughout this section a step goes from layer $j$ to layer $j-1$.
Write $l=L$ for a regular step and $l=L_*=5L$ for an inclined
first step.  Put $t=t_j$, $t'=t_{j-1}$, $p=p_j$,
$M=\mathfrak m_j$, $M'=\mathfrak m_{j-1}$, and
$g=b^{-1/2}$, where $b\in[H_j/2,2H_j]$.  Thus $t$ is comparable to
$gp$.  All geometric cutoffs use the fixed reference scales $H_j,t_j$,
including when two precise couplings are compared.  A history records
the linear observations already performed.  Histories may start with
one inclined observation and then use regular observations.

\subsection{The retained class and the step theorem}

We use the retained densities and norms of Section~\ref{sec:setup}.  The tangent coordinates at an anchor $o$ are
\[
 y_o(x)=\operatorname{proj}_{q_o^\perp}q_x.
\]
Coefficient tensors are invariant under the full orthogonal group of
$q_o^\perp$.  Their value is consequently independent of an
orthonormal basis of that plane.  The odd canonical slots $P_5,I_3$
vanish; it is convenient to retain their zero entries in the canonical
vector $\mathbf P$.  Quadratic packets are averaged over spatial
quarter turns and compensated to have zero affine Hessian.  The two
factors of a quadratic polynomial are symmetrized before compensation.
Indeed an $O(2)$-invariant color contraction is the dot product, and
quarter-turn averaging of its symmetric spatial quadratic form gives
a scalar multiple of the identity.

The same quarter-turn convention is imposed on all histories.  Every
choice also commutes with simultaneous reflection of the geometry and
the input data, as in Section~\ref{sec:free}; a precision lacking a
reflection symmetry is not itself averaged over that reflection.
Paths can be chosen by retaining the finitely many coordinate-order
paths together, or by taking their union.  A connected expression keeps
its individual joining paths before its output mask is enlarged.
Each bulk regular-error packet is invariant under simultaneous spin
transformations and under spatial inversion about its anchor.  Its
mask is strengthened to a common invariant graph.  As in
\Rref{rg:class}, the exact off-mask correction is retained in the
covering gas.

The error norm uses actual classical derivatives on the open graph
domain, with $k=8$.  At a spin $q_x$, the rotation directions are
$v_x\times q_x$, with the spatial weights specified in the common
setup.  At a single site they give uniformly equivalent tangent norms, and
fixed small spherical charts have uniformly bounded coordinate
derivatives.  For graph norms, chart chain rules also multiply spatial
weights; the resulting mesoscopic and load factors are tracked in
Lemma~\ref{lem:rg-error-contraction}.  No smoothness is required of an error away from its graph domain.
Integration factors with Gaussian color indices retain their full
contracted tensor.  Thus internal invariance holds for each integrated
packet, although it need not hold for a component before its indices
are contracted.  The same convention applies to factors carrying
external vector indices in Section~\ref{sec:sources}.

On a torus, a calculation is declared winding when its complete lifted
record crosses the prescribed shortest-period threshold.  Short orbit
calculations use the bulk symmetry of the lifted labels, even when the
torus itself is not invariant under those symmetries.  Winding terms
retain their winding price and need not have the bulk normalization.
The initial nearest-neighbor representation has these conventions:
all witnesses of a deleted masked row are joined in its component.
Fixed enlargements of graphs, including projection to an inclined
coarse grid, change only the constant in the load-projection estimate.
The short-record threshold can be decreased and the minimum admitted
period increased to accommodate them.

\begin{theorem}[Exact step and comparison]
\label{thm:rg-step}
Fix $\upsilon>0$ sufficiently small.  There are choices of $L$, the
canonical caps and comparison weights, $P_0$, and finally $H$ such that
integration against one normalized block observation sends every
admitted retained density at $b\in[H_j/2,2H_j]$ to an exact
representation of the same form in layer $j-1$, with all renewed
shape caps.  The representation may be iterated when its new coupling
belongs to $[H_{j-1}/2,2H_{j-1}]$; that admission is established
separately by shooting in Section~\ref{sec:shooting}.  Its precise
coupling is
\begin{equation}
 b'=b+\alpha_h(\mathbf P)+\frac{\kappa_h(\mathbf P)}b+\Delta,
 \qquad |\Delta|\le C_LH_j^{-1.05}.
 \label{eq:orig-4}
\end{equation}
Here $h$ is the free history, and the bulk scalar, canonical slots,
and kicks use the prescriptions of \Rref{rg:canonical-map} adapted
below to $S^2$.  All retained norm caps are renewed.

For two regular steps with common geometric scales, let
$\lambda=b_2-b_1$, and let $u$ be the fixed weighted shape discrepancy:
canonical coefficient differences, regular-error difference divided by
its cap, and covering-gas difference divided by its cap.  The regular
error difference is measured through order $k-1$.  Then
\begin{align}
 u'&\le q u+C_L\operatorname{poly}(H_j)\epsilon_h
       +C_L(\log H_j)^C|\lambda|/H_j,\nonumber\\
 |\lambda'-\lambda|&\le C_Lu+C_L\operatorname{poly}(H_j)\epsilon_h
       +C_L(\log H_j)^C|\lambda|/H_j,
 \qquad q<L^{-2+\upsilon}.
 \label{eq:orig-5}
\end{align}
For identical histories $\epsilon_h=0$.  If the histories share their
last $r$ regular observations, one may take
$\epsilon_h=(A_0/L^2)^r$, with the constant $A_0$ of
Section~\ref{sec:free}.  The same-step Lipschitz bounds with no history
change hold also for an inclined first step.  The representation classes and constants
are common to all the admitted histories and periods.  Signed covering
activities are allowed: the logarithm used in constructing the regular
exponent is the convergent connected logarithm of the small background
factor, based at one.  The mandatory-cover factor is retained as a sum,
and no logarithm of the full signed density is taken.
\end{theorem}

We prove the theorem in five stages: a sector decomposition isolates
large differences; ambient Gaussian completion treats the remaining
spins; joint estimates control products of the resulting factors;
canonical extraction identifies the finite Taylor terms; and
normalization returns the remainder to the retained class.  The last
stage uses the strengthened contraction proved below, rather than the
$L^{-1}$ contraction of \Rref{rg:error-contraction}.

\subsection{Sector geometry and the energy identity}
\label{subsec:rg-sectors}

We first specify the geometric data that remain fixed during each
Gaussian integral.  All distances are measured in fine-lattice units,
with a coarse position embedded at its block anchor.  Primitive free
kernels are truncated at a radius $c_LM$, and inverse windows are
chosen on the same scale.  The constant is reduced so that the sum of
radii in each of the finitely many local compositions, including block
diameters, reference displacements and kinetic-mask reads, is less
than $M/100$.  This is possible because $L$ is fixed and the mask
radius is $O_L(\log H_j)=o(M)$.  Omitted paths have weighted sums
$C_Le^{-c_LM}$ and are kept as separate factors.  Arbitrarily long
inverse and determinant chains keep all their successive windows and
endpoints in their records.

A smooth edge profile is one for $|d_eq|\le t/4$ and zero for
$|d_eq|\ge t/2$.  A smooth observation profile is one for
$|V_Y-\widehat m_Y|\le Wt$ and zero for
$|V_Y-\widehat m_Y|\ge2Wt$, with $W$ fixed sufficiently large.
Expanding each profile plus its complement gives an exact partition.
A chosen complement is called an input or observation \emph{seed}.
Every true output bad edge is also a mandatory seed.  Join seeds whose
radius-$50M$ neighborhoods intersect.  A resulting component, together
with its assigned factors, is called a \emph{core}.

The assignments are those of \Rref{app:factor-ownership}: spins within
$3M$ of the seeds are frozen and retain Haar integration; numerical
energy rows within $5M$ are assigned to the core; stationary squares
within $8M$ and default profiles within $14M$ have the same owner.
Leading determinant changes have an owner within $3M$ plus the inverse
window radius.  The larger collections do not duplicate the numerical
rows.  Distinct $50M$ footprints are disjoint, so every assignment is
unique.  Maximality of a component is imposed by compatibility of its
footprints in the pattern sum, not by an extra distant predicate in a
single core factor.  A \emph{complete support} records every numerical
argument, mask, eligibility test, present or absent status query,
window and joining path used by the factor.  This convention is
essential for factorization later.

Here are the energy rows whose ownership has just been fixed.  Let
$d,d'$ denote fine and coarse differences, let $\ell_s,\ell'_s$ be the
truncated free edge operators, and put
\[
 X=(\sqrt{c_0}\,dq,\ell_s dq),\qquad
 Y=\sqrt a(V-\widehat m),\qquad
 X'=(\sqrt{c_0}\,d'V,\ell'_s d'V),\qquad a=1.
\]
Write $X_2,X'_2$ for the second components.  For a vector $z$,
$\operatorname{clip}_t z$ has length $\min(|z|,t)$ and the same
direction, with value zero at $z=0$.  With the output kinetic mask $m'$,
define
\[
 \zeta=(\sqrt{c_0}\operatorname{clip}_{t'}d'V,
                         m'\ell'_s d'V),
 \qquad (\nu,w)=\mathcal M_s\zeta.
\]
The free bounds imply
$|\zeta|+|w|\le Ct$ and $|\nu|\le Ct/L$ entrywise.  The input kinetic
energy plus observation energy minus output kinetic energy equals the
sum of
\begin{align*}
 \mathcal F_e&=S_t(r_e)+\tfrac12m_e|X_{2,e}|^2
                      -\nu_e\cdot X_e+\tfrac12|\nu_e|^2,&
 \mathcal B_Y&=\tfrac12|Y_Y-w_Y|^2,\\
 \mathcal C_E&=\zeta_E\cdot X'_E-\tfrac12|\zeta_E|^2
           -S_{t'}(|d'_EV|)-\tfrac12m'_E|X'_{2,E}|^2,\\
 \mathcal U_E&=\zeta_E\cdot[\mathcal M_s^*(X,Y)-X']_E,&
 \mathcal N_i&=\tfrac12| (\mathcal N_s\zeta)_i|^2,\\
 \mathcal Z_E&=\tfrac12\zeta_E\cdot
       [(\Id-\mathcal M_s^*\mathcal M_s-\mathcal N_s^*\mathcal N_s)
                                                     \zeta]_E.
\end{align*}
This is the identity \Rref{rg:ledger}, applied to real ambient fields.
The cross terms cancel because $(\nu,w)=\mathcal M_s\zeta$; the last
row accounts exactly for truncation of the free isometry.  In
particular \Rref{rg:ledger-errors} gives
\[
 |\mathcal Z_E|\le C_Le^{-c_LM}t^2,
 \qquad
 |\mathcal U_E|\le C_Lte^{-c_LM}
       +Ct\sum_Ye^{-c\operatorname{dist}_f(E,Y)/l}|m_Y-\widehat m_Y|.
\]
Here $\operatorname{dist}_f$ denotes the fine-coordinate distance
between the embedded coarse sites; dividing by $l$ measures the kernel
decay in coarse units.  Thus its sum over $Y$ is bounded independently
of $L$.  These formulas define the rows before any small-angle
expansion.

Select the $\mathcal F,\mathcal B$ rows farther than $2.5M$ from every
seed.  Each row incident on a nonfrozen column is selected.  The
retained defaults, including their closed derivative supports, imply
that all spins read by a selected assigned row are within $C_Lt$ of
their block reference.  For regular cells this is the coordinate-path
proof of \Rref{app:principal-log}; for inclined cells it uses the paths
and mean estimate of Section~\ref{sec:free}.

The support prescriptions of \Rref{app:rg-geometry} are retained in
full.  In particular an optional factor $e^V$ with merely measurable
spin dependence is prepared as
$1+\chi_{\mathrm{att}}(e^V-1)$, not by opening
$\chi_{\mathrm{att}}e^V$.  Here $\chi_{\mathrm{att}}$ is the product
of its attached chart protectors and alignment profiles.  It is one
on the original sector, remains attached to the selected letter, and
does not create an order-one failure outside its plateau.  A separate
global chart profile removes nonprincipal exponential preimages.

For a calculation with empty output inventory, all queried output
statuses are fixed to no flag, clipping is replaced by the identity,
and good output masks by one.  The resulting analytic formula is
extended through graph gradients $<4t'$ and cut off only on a larger
graph domain.  It agrees with the physical formula on the outgoing
regular mask.  If output inventory is nonempty, its statuses stay
fixed and no field derivative is taken.  In neither case is a
discontinuous output predicate differentiated across its jump.

\subsection{Ambient Gaussian completion}
\label{subsec:rg-ambient}

For a nonfrozen site put $T_x=V_{[x]}$ and use the principal intrinsic
coordinate
\[
 u_x=\operatorname{Log}_{T_x}q_x\in T_x^\perp,
 \qquad q_x=\operatorname{Exp}_{T_x}u_x.
\]
The sector bounds give $|u_x|\le C_Lt$.  Frozen coordinates read by
selected rows use the same logarithm multiplied by a smooth radial
cutoff inside its injectivity radius $\pi$, and extended by zero.
This is a globally smooth function of the pair of spins.  Protectors
are supported strictly inside that chart and equal one on the
required plateaus.  Unlike a tangent frame, this construction is
intrinsic.

\begin{lemma}[A normalized ambient completion]
\label{lem:rg-ambient}
Fix a primary pattern, the frozen spins, the tags in the observation,
and the output data with their statuses.  Let $D$ be the scalar column
matrix, on the nonfrozen sites, of the selected rows of
\[
 D_0=(\mathsf A_s d,-\sqrt a Q),\qquad A=D^*D.
\]
Adjoining the normalized scalar Gaussian with negative log density
$b|Dv|^2/2$ and setting
\[
 \xi_x=u_x+T_xv_x
\]
gives an exact integration in $\R^3$ at every nonfrozen site, with
Lebesgue Jacobian one for this change of variables.  The matrix $A$
satisfies $c_L\Id\le A\le C_L\Id$, uniformly in the pattern and
period.  No cutoff on $v$ is required.

For a selected row with reference $T_i=V_{o(i)}$, let its affine
constant be the projection onto $T_i^\perp$ of the physical row
$R_i(0,V)$, where $R=(X-\nu,Y-w)$.  Insert the smoothly extended
frozen coordinates in $D_0$ and denote the resulting affine row by
$D_i\xi+F_i$.  On the selected-row and smooth angular-cutoff supports,
the additional error caused by the scalar coordinate in subtracting
the affine square is bounded in absolute value by
\begin{equation}
 C_LM^D\left(t^2\sum_x|v_x|+t\sum_x|v_x|^2\right).
 \label{eq:orig-6}
\end{equation}
The sums run over that row's complete coordinate reads.  The same
estimate holds on the enlarged no-output-flag graph domains.
\end{lemma}

\begin{proof}
Because every row incident on a nonfrozen column is selected, $A$ is
the corresponding principal restriction of
$d^*\mathsf A_s^*\mathsf A_s d+aQ^*Q$.  Block Poincar\'e, with vectors
extended by zero into frozen sites, gives
\[
 \norm{v}_2^2\le C_LL^2\bigl(\norm{dv}_2^2+\norm{Qv}_2^2\bigr).
\]
The same argument uses the cell paths for an inclined step.  Together
with the fixed positive direct-edge coefficient it proves the stated
ellipticity, as in the paragraph preceding \Rref{rg:window-bounds}.
The empty exterior causes no exception: the Gaussian integral is then
one.

At fixed $T_x$, the orthogonal sum
$T_x^\perp\oplus\R T_x=\R^3$ is an isometry.  Hence $(u_x,v_x)\mapsto
\xi_x$ has Jacobian one.  The original spin measure contributes only the dimension-two spherical
Jacobian
\[
 J(u)=\frac1{4\pi}\frac{\sin|u|}{|u|},\qquad J(0)=\frac1{4\pi},
\]
where the value at zero is taken by continuity.  The scalar integral
has been normalized, so no degree
of freedom has been added to the physical measure.

All reference spins in a selected row's reads differ by at most
$C_LM^Dt$, while physical angles are at most $C_Lt$.  At zero angle,
the normal component of $R_i(0,V)$ is
$O(C_LM^D(t^2+e^{-c_LM}t))$.  Taylor expansion of the stack gradients
and chord mean proves this bound; scalar kernels commute with a
constant reference rotation.  The linear variation of the physical
row is $D_0u$, with remainder $O(C_LM^Dt^2)$.  Meanwhile
\[
 D_0(Tv)-T_iDv=O(C_LM^Dt)\max_x|v_x|.
\]
The product of $T_iDv$ and the leading physical tangent row is zero.
The remaining cross products give the first term of
Equation~\eqref{eq:orig-6}, and the scalar-square mismatch gives its
second term.  Arbitrarily accurate truncation errors are included.
For a symmetry tie between finitely many references we use the average
of the projected affine row and retain every reference read.  The
same estimate can be made with any one of those references.

The enlarged no-flag domains merely replace the fixed small-angle
constant by a larger one.  No small-angle bound has been asserted for
a nonselected core row, whose physical energy remains unexpanded.
\end{proof}

We next make Gaussian localization explicit.  For a nonfrozen site $x$
let $\Omega_x$ be its inverse window, restricted to nonfrozen sites.
For $a_0\ge0$ define the column-window inverse and its residual by
\[
 (\Pi_{a_0})_{yx}=\mathbf1_{y\in\Omega_x}
               (A_{\Omega_x}+a_0)^{-1}_{yx},\qquad
 E_{a_0}=(A+a_0)\Pi_{a_0}-\Id.
\]
We use $a_0$ here to distinguish the resolvent parameter from the
observation precision $a=1$.  The estimates \Rref{rg:window-bounds}
give weighted row and column bounds $C_L/(1+a_0)$ for the inverses
and $C_Le^{-c_LM}/(1+a_0)$ for the residuals, including their history
differences.  Define
\[
 C_s=\tfrac12(\Pi_0+\Pi_0^*),\qquad
 \mu=-\Pi_0^*D^*F,\qquad r_s=D\mu+F.
\]
For large $H$, $C_s$ and every principal marginal have fixed-$L$
upper and lower spectral bounds.  Use one common ambient Gaussian
$Z\sim N(0,C_s\otimes\Id_{\R^3})$ and substitute $\xi=\mu+gZ$.
Expanding the affine squares gives the exact density-ratio identity
\Rref{rg:gaussian-ratio}; its remaining exponent is
\[
 \tfrac12Z^*(C_s^{-1}-A)Z
             -g^{-1}(A\mu+D^*F)^*Z.
\]
With $E_s=AC_s-\Id$, the corrections are the convergent series
\[
 C_s^{-1}-A=\sum_{r\ge1}(-E_s)^rA,
 \qquad
 \log\det C_s=-\Tr\log A+
       \sum_{r\ge1}\frac{(-1)^{r+1}}r\Tr E_s^r,
\]
together with the resolvent expansion for $\log A$ in
\Rref{rg:ratio-series}.  Each residual path contains at least one
exponentially small residual hop.  For every fixed required support
exponent its summed bound is the chain bound \Rref{rg:chain-price},
with the same bound for one history difference.

The scalar normalization in Lemma~\ref{lem:rg-ambient} cancels one
color in the leading determinant.  Its coefficient is therefore two,
although the Gaussian used for all estimates has three ambient
colors.  In the explicit local determinant compensation of
\Rref{app:explicit-core}, replace the coefficient $3/2$ by $1$ and
use $\kappa_g=(2\pi g^2)J(0)$.  For example the leading window scalar
at $x$ is
\[
 h_x=-\int_0^\infty
    \bigl((1+a_0)^{-1}-(A_{\Omega_x}+a_0)^{-1}_{xx}\bigr)\,\dd a_0.
\]
Its empty-pattern value retains the site's phase modulo the block
lattice.  The frozen-site compensations and changes near a hole are
assigned as above.  Residual scalar-determinant terms are the same
exceptional paths with fixed coefficients.  Thus their complete
supports and locality are unchanged.

\begin{lemma}[Prediction of a spin]
\label{lem:rg-prediction}
On a stencil with the margins of \Rref{app:read-sets}, in the empty
pattern or away from nearby cores, put
\[
 q_x^0=\operatorname{Exp}_{T_x}
                     (\operatorname{proj}_{T_x^\perp}\mu_x).
\]
Then the three estimates \Rref{rg:prediction} hold with this spin:
\[
 r_s=O\bigl(C_LM^D(t^2+te^{-c_LM})\bigr),\qquad
 |dq^0|\le Ct/L+C_LM^Dt^2,\qquad
 |\mu|\le Ct+C_LM^Dt^2.
\]
They hold with the fixed normalized derivatives and with a single
history or precise-coupling difference.  The leading constants are
uniform before $L$ is chosen.  Everywhere, including on extensions
near cores, $|\mu|\le C_LM^Dt$.
\end{lemma}

\begin{proof}
Use a tangent chart at a coarse reference spin $e_0$ and denote the
first tangent data of $V$ by $w_0$.  At angle zero the fine spin
prescription has tangent $w_0([x])$.  Its variation in $u$ adds
$D_0u$ to the rows.  With full kernels, the choice
\[
 u(x)=(P_hw_0)(x)-w_0([x])
\]
solves the linearized affine equations with zero residual, by the
ambient and harmonic free identities.  The normalized mean has
tangent differential the identity.  Projecting the constant row does
not alter its first tangent.  Invertibility of $D^*D$ therefore makes
this the exact first tangent of the affine Gaussian mean, with zero
normal component.  Only linear data are used in this identity.

Window and kernel truncation change the comparison by
$C_Le^{-c_LM}$ in all required weighted kernel moments.  On the
supplied output graph the stencil lies in an $O_L(Mt)$ chart.  Missing
positions outside it may be filled by zero in that chart; the retained
margins make their contribution exponentially small.  Based at a
block spin, output chord bounds give coordinate size
$Ct(1+|z-[x]|)$.  The summable weighted rows of $P_h$ therefore give
$Ct/L$ for a one-fine-edge difference, and $Ct$ for the mean with its
constant part subtracted.  Products and remainders of the smooth spin
and projection maps cost $C_LM^D$ times the indicated products of
sizes, yielding the quadratic errors above.

Each normalized derivative contributes $t$ times the order-eight
spatial direction weight.  The leading row moments sum that weight;
other contributions cost a fixed polynomial in the stencil size.
Changing the direction anchor inserts the corresponding distance
weight in each slot, paid by the complete record and its load.  Every
map derivative beyond first order has at least two factors of $t$.
Stationary rows have zero first tangent except for the exponentially
small window discrepancy.  This also proves the needed mean and
stationary-row differences by the free-kernel comparison bounds.
Frozen-angle extensions outside these local prediction regions only
need the fixed polynomial bounds.  A cutoff of their defining data
on a larger graph domain gives the global bound for $\mu$, without
requiring a global spherical logarithm.  This verifies, in particular,
the empty-pattern predictions at every center used in a mask transfer.
\end{proof}

\subsection{Products of factors and large-gradient reserves}
\label{subsec:rg-products}

The completion has reduced each sector to one Gaussian expectation,
frozen Haar integrals and finite tag averages.  Factors in this common
expectation are generally correlated.  We establish product estimates
before expanding their connected contributions.

For a core $c$, let $s_c$ be its number of seeds and $I_c$ its Gaussian
read set.  In its energy sum subtract each selected assigned affine
square exactly once, and include the assigned stationary squares,
profiles, local inventories, attached protectors and normalization
compensations.  Denote the resulting factor by $B_c$.  Thus it is the
factor \Rref{app:explicit-core}, with the scalar modification in
Lemma~\ref{lem:rg-ambient}.  An input primary that belongs to a true
bad-bond inventory is supplied once by its old covering label; every
other selected input primary retains $\mathbf1_{r_e\le t}$.  Nonprimary
edges already satisfy $r_e\le t/2$.  These local inventory rules are
exactly the original covering constraint.

\begin{lemma}[Ambient core and tail estimates]
\label{lem:rg-reserves}
The core factor retains the reserve $e^{-c_Lp^2s_c}$ of
\Rref{app:explicit-core}, apart from the additional envelope
\begin{equation}
 \exp\left\{e_0\left(M^Ds_c+
                         \sum_{x\in I_c}|Z_x|^2\right)\right\},
 \qquad e_0=C_Lg\operatorname{poly}(M,p).
 \label{eq:orig-7}
\end{equation}
The bounds hold on every attached profile's closed derivative support.
On a selected unassigned row the exponent to be opened has, on its
angular-cutoff support, the envelope
\begin{equation}
 C_Lg\operatorname{poly}(M,p)
                 \left(1+\sum_{x\in I_i}|Z_x|^2\right),
 \label{eq:orig-8}
\end{equation}
where $I_i$ is its Gaussian read set.  Such a row is the sum of an
ordinary letter supported where $\max_{I_i}|Z_x|\le2p$ and a tail
letter requiring $\max_{I_i}|Z_x|\ge p$.  The former has ordinary
majorant $C_Lg\operatorname{poly}(M,p)$; the latter has an exceptional
Gaussian reserve.  Both assertions hold with every fixed required
normalized derivative and a single marked discrepancy.
\end{lemma}

\begin{proof}
The charging proof following \Rref{app:explicit-core} uses direct
positive squares, failed-mask witnesses, output reserves, and the
mean-error bound.  These estimates are independent of the dimension
of the sphere.  Lemma~\ref{lem:free-inclined-geometry} supplies its mean bound,
and Corollary~\ref{cor:free-small-chord-mean} supplies the stronger
$C_LT_j^2$ bound when all relevant chords are at most $T_j$.  The numerical-square comparisons are the ones in
Lemma~\ref{lem:rg-ambient}.  Consequently the lower bound for the core
energy, including its profile derivative supports, is unchanged
except for Equation~\eqref{eq:orig-6} summed over selected assigned
rows.  No direct square is charged twice.

Now $v_x=T_x\cdot(\mu_x+gZ_x)$, and
$|\mu_x|\le C_LM^Dt$.  Multiplication of
Equation~\eqref{eq:orig-6} by $b=g^{-2}$, with $t\asymp gp$, bounds
the added exponent by a fixed polynomial times
$g(1+\sum|Z_x|^2)$.  Row counts and local overlaps are polynomial in
$M$ and independent of $s_c$ at a fixed site.  Incorporating them into
$e_0$ gives Equation~\eqref{eq:orig-7}, and the single-row version is
Equation~\eqref{eq:orig-8}.  The remaining core factors have the bounds
in \Rref{app:explicit-core}, including $e^{-c_Lp^2s_c}$.

For an unassigned row insert smooth products of radial profiles in
$Z$, equal to one through $p$ and zero beyond $2p$.  On the ordinary
part the exponent and its fixed derivatives are
$C_Lg\operatorname{poly}(M,p)$.  A derivative in a reference spin
either preserves its small difference from the other references or
replaces that difference by $t$ times a direction weight.
Differentiating $v$ likewise inserts only a fixed polynomial on these
supports.  The tail part has the same exponential envelope as a core
and forces a Gaussian coordinate of size at least $p$.  Derivatives
of its profiles have polynomial cost.  This splitting changes no
geometric statuses or row ownership.

All extended chart factors agree with their physical values on the
original sector; the global profile suppresses other preimages, and
the protectors remain attached.  Thus the estimate has not enlarged
the integration by including an uncontrolled branch.  The marked
discrepancy assertion is proved below with the product comparison.
\end{proof}

For clarity, we specify the remaining factors to which the joint
estimate will be applied.  This is the eight-class preparation of
\Rref{rg:prepared}, with the row split just described.
\begin{enumerate}[label=(\roman*)]
\item Selected unassigned row corrections and the smooth mean-defect
row $\mathcal U$ are opened with angular plateaus containing their
full-sector values.  Their ordinary order is $g$.
\item Unassigned stationary squares have order
$g^2\operatorname{poly}(M,p)$.  Unassigned $\mathcal C$ rows vanish;
$\mathcal N,\mathcal Z$ rows have exponentially accurate exceptional
bounds, using the free null identity.
\item An old canonical polynomial or regular error is called remote
when its entire projected graph is farther than $14M$ from all seeds.
There its graph mask is replaced by a smooth product equal to one
through $t/2$ and zero from $2t$.  The error is defined through $4t$,
so the resulting function is smooth.  A nonremote term retains its
actual mask and attached protectors and meets a core.  Path estimates
give \Rref{rg:old-polynomial-bound}; old errors retain their actual
norm caps.  A common chart on an arbitrarily large graph is not used.
\item Kinetic truncation tails retain both paths and all mask reads.
They have arbitrary-power accuracy with exponential path weights.
Their input evaluations use the same smooth or protected prescription;
output evaluations use the fixed-status convention.
\item Inverse and determinant chains are kept literally.  A ratio
exponent $Q_\alpha$ has at most two Gaussian endpoints and satisfies
$|Q_\alpha|\le e_\alpha(1+\sum_{I_\alpha}|Z_x|^2)$.  Even the eighth
roots of $e_\alpha$ are summable with every fixed support exponent
required here, by \Rref{rg:chain-price}.
\item The spherical Jacobian is opened as a smoothly cut-off ratio
$J(u)/J(0)-1$.  It starts quadratically and has order
$g^2\operatorname{poly}(M,p)$.
\item Remaining edge, observation and global-chart defaults are
opened as one plus signed failures.  In a remote observation default
the mean is normalized only on the smooth branch bounded away from
zero.  Its failure and all its nonzero derivative supports force a
large Gaussian coordinate, by Lemma~\ref{lem:rg-prediction}.  Projecting
$\xi$ to the tangent plane cannot invalidate that implication.
\item An old covering weight keeps its supremum envelope and its
protectors.  Its nonempty inventory meets a core.  Its merely
measurable physical-spin function is never differentiated.
\end{enumerate}
The three-color Gaussian scalar powers, the corrected two-color
leading determinant, the Haar constant and the observation denominator
have already been extracted.  There is no ordinary factor of order
zero left unaccounted for.

\begin{proposition}[Joint integration and comparison]
\label{prop:rg-joint}
For every finite admissible selection $\mathcal A$ of prepared
factors in a fixed pattern, including the compulsory cores, there
are nonnegative majorants $a_\alpha$ such that
\[
 \left|D^r\E\prod_{\alpha\in\mathcal A}F_\alpha\right|
 \le C_k\left(1+\sum_{\alpha\in\mathcal A}s_\alpha\right)^{d_k}
                        \prod_{\alpha\in\mathcal A}a_\alpha,
 \qquad r\le k.
\]
The expectation includes the Gaussian, frozen spins and tags.
For nonempty output inventory this asserts only the undifferentiated
and parameter-comparison estimates.  A single marked discrepancy
retains its size in this product bound.  An ordinary primitive of
order $r_0>0$ has summed majorant
$C_Lg^{r_0}\operatorname{poly}(M,p)$ with the prescribed support
weights.  Exceptional factors have arbitrary fixed power accuracy,
and seeded connected sums retain an $e^{-p^{3/10}}$ reserve after
all decorations.  Disjoint complete supports factor exactly.
\end{proposition}

\begin{proof}
The unmodified factors have the concrete envelopes of
\Rref{supp:joint-majorants}.  Their proofs use the prediction lemma,
path bounds, and the graph $C^k$ domains just verified.  Rotation
fields realize tangent derivatives with uniformly smooth bounded
coefficients.  Projection derivatives outside a smaller plateau
insert fixed Gaussian polynomials, already allowed by those estimates.

Include the new tail-row failures in the simultaneous-failure bound
\Rref{rg:joint-rarity}.  Distinct rows have polynomially many overlaps
and each forces a coordinate at least $c_Lp$.  Their Gaussian rarity
therefore has the same form as the previous failures.  For a product
of the extra envelopes in Equations~\eqref{eq:orig-7} and
\eqref{eq:orig-8}, core components are disjoint and the number of
row-tail reads at one coordinate is $O_L(M^D)$.  The diagonal charge
in the combined quadratic exponential has operator norm
$C_Lg\operatorname{poly}(M,p)=o_H(1)$.  Its trace and constant terms
cost at most $C_Le_0M^{D'}$ times the sum of core loads and tail-row
counts.  The determinant estimate \Rref{supp:gaussian-product}, at
any fixed larger moment exponent, thus applies.  Increasing the
fixed H\"older exponent if necessary, these costs are paid by the
exceptional reserves.  The summable endpoint charges of ratio letters
and the deterministic ordinary envelopes give precisely
\Rref{rg:joint-product}.

Complete supports include the windows determining $C_s$ and all
status reads.  Its finite range then gives zero covariance between
Gaussian read sets whose complete supports are disjoint.  The frozen
Haar and tag measures are product measures.  Hence
\Rref{app:local-marginals} applies to \emph{marginals}; it does not
assert independence of conditional Gaussians on overlapping supports.

It remains to justify derivatives of hard spin evaluations and the
marked comparison.  These are addressed in the next two paragraphs
as part of the proof, because separate pointwise Taylor accuracy
would not imply the asserted joint estimate.

\emph{Transport at a hard read.}
At fixed $T_x$ the map
\[
 \xi_x\longmapsto
 \left(\operatorname{Exp}_{T_x}
             (\operatorname{proj}_{T_x^\perp}\xi_x),
                              T_x\cdot\xi_x\right)
\]
has inverse
$\operatorname{Log}_{T_x}q_x+T_xv_x$ on the protector chart times
the entire real normal line.  Its fixed chart derivatives have only
polynomial growth in $|v_x|$.  Holding $(q_x,v_x)$ fixed while a
parameter varies is achieved by the Gaussian transport
\[
 W_x=g^{-1}\left\{
    \partial(\operatorname{Log}_{T_x}q_x+T_xv_x)
                 -\partial\mu_x-(\partial g)Z_x\right\}.
\]
Here $\partial\mu$ is evaluated at fixed frozen variables and tags;
output spins may move only in the no-flag calculations.  Extend this
vector field by a smooth cutoff beyond the protector support but
strictly inside the chart.  It gives the transport and divergence
bounds of \Rref{rg:transport}, with the allowed powers of $g^{-1}$
and fixed polynomials in $Z,M,p$ and directional distances.
Different sites can use different charts simultaneously.

For merely measurable input functions the same assertion follows by
changing variables to $(q_x,v_x)$ on the protected reads, with
protectors attached, and differentiating while those arguments are
fixed.  This is the change-of-variables version of integration by
parts.  At every finite cutoff and for every finite product, the
positive quadratic charges above remain strictly below the Gaussian
quadratic decay.  The chart changes have fixed margins; the unbounded
normal coordinate changes orthogonally, with a scaling and shift.
Differentiation under these integrals is therefore justified by
Gaussian domination.  Pullback to the common Gaussian marginal gives
the same score estimates.  The inverse covariance on any union of
reads has norm $O_L(1)$ by ellipticity, so its density scores have
fixed polynomial moment costs as in \Rref{supp:product-reserve}.
No discontinuous predicate in a moving angle is tested.

\emph{One parameter or history difference.}
Fix the geometry, output data, tags and frozen arguments, enlarge
supports to contain both endpoints, and use the same cutoffs.
Interpolate $C_s$, $g$ and $\mu$ linearly.  Evaluate $\xi,u,v,q$ at
these common interpolated arguments.  Pure window data, determinant
chains, Gaussian-ratio exponents and their recorded paths can use the
linear interpolation of their endpoint formulas at common $Z$.
For residual row and core exponents interpolate \emph{after} each
endpoint construction: regard them as functions of common
$(\xi,V)$ and fixed statuses, and interpolate those functions.
No energy identity is assumed for a mixture of endpoint kernels.

Both endpoint core budgets hold on these common physical/chart tests.
Throughout the interpolation, $|v|/g$ is bounded by $\max|Z|$ plus a
fixed polynomial in $M,p$.  The first tangent prediction of the
interpolated mean is the convex combination of the endpoint first
tangents.  Passing to its spin adds only
$O_L(M^Dt^2)$ and arbitrarily accurate window errors.  Thus the
small-$Z$ contradiction that proves every failure implication still
works, and the global chart failure uses the same bound on $\mu$.
All H\"older estimates therefore hold along the interpolation.

Let $\varepsilon=\epsilon_h+|\lambda|/H_j$.  On aligned supports,
\[
 |\mu_2-\mu_1|\le C_LM^Dt\varepsilon,
\]
with its normalized smooth derivatives.  The corresponding stationary
rows have the same discrepancy factor because their first tangent
vanishes up to window accuracy.  The projected affine maps have
bounded cutoff extensions and exponentially summed windows, so these
bounds remain valid with the polynomial hard-transport losses.  In a
row-square difference, the leading physical term $D_0u$ cancels
separately at each endpoint; its remainder difference has the same
second-order small factors times $\varepsilon$.  In
$D_0(Tv)-T_iDv$, a reference difference, or its derivative, multiplies
the kernel discrepancy; frozen normal columns are zero.  The mean
defect in $\mathcal U$ is already quadratic, and its full-kernel
cancellation holds separately at the two endpoints.  Kernel-tail
differences retain their exponential factor.

Consequently an ordinary letter's parameter derivative is bounded by
its majorant times $\varepsilon$, with fixed polynomial losses in
$p,M,1+s_{\mathrm{tot}}$ and fixed Gaussian polynomials.  A normalized
direction supplies $t$ times its spatial weight; the $t^{-1}$ cost of
an angular cutoff is paired with that factor or with the mean/window
discrepancy.  Unbounded tail letters retain the same discrepancy in
the exceptional envelope and polynomial scores.  Hard uses retain it
by the displayed inverse transport, attached profiles and endpoint
exponents evaluated at common physical arguments.  An input-list
difference, handled separately by zero padding if needed, keeps its
supremum difference at a measurable read.

Only a fixed number of scores or derivative slots occur.  At fixed
H\"older order their Gaussian moments grow polynomially in the
number of reads, hence in $M$ and the total load, even when reads
overlap.  This proves the marked discrepancy estimate with $k-1$
additional derivatives on a regular output.  A derivative placed on
a changed regular input uses at most $k-1$ of that difference; a map
change acting on an unchanged input uses at most its $k$ derivatives.
Ordinary derivative costs have remained of ordinary order.  The
larger hard-transport costs are paid only by exceptional reserves.

Finally row shifts, chart reads, field rank and the new profile types
have bounded stencil complexity and polynomial overlap.  Count and
site-hit bounds therefore use the same loads, with enlarged fixed
polynomial exponents.  Their spare support exponents absorb the
load polynomial in the displayed product estimate.  The order and
seeded bounds now follow from the same summations as
\Rref{rg:joint}; this completes the proof, including the discrepancy
part of Lemma~\ref{lem:rg-reserves}.
\end{proof}

The hypotheses of the exact connected reassembly
\Rref{rg:connected} are now available: joint products, site-hit bounds,
the tree smallness condition, and complete support and inventory
rules.  Its Mayer expansion is used only for the background without
mandatory output flags.  Components with flags retain their complete
covering weights.  We obtain \Rref{rg:preliminary-output}, with the
exact mask-strengthening corrections described there.  For countable
factor lists, first perform the algebra on finite selections; the
volume-exponential absolute majorant from the joint and tree bounds
permits passage to the full expression, including its derivatives.
Positivity of a truncated signed gas is not needed.  Thus the output
representation accounts for all mandatory covers, not only the
sector with no seeds.

\subsection{Canonical extraction and its diagonal contraction}
\label{subsec:rg-canonical}

The connected representation just obtained is exact.  We next separate
its finite-order Taylor polynomial from a smaller regular remainder.
The coefficient calculation uses full bulk kernels and aligned charts,
so it is independent of periods, masks and chart cutoffs.

In the empty pattern let $R=(X-\nu,Y-w)$ with
$(\nu,w)=\mathcal M X'$ and full kernels.  The mean-normalization
defect is
\[
 U_E=X'_E\cdot\bigl[\mathcal M^*
                    (0,\sqrt a(m-m/|m|))\bigr]_E.
\]
Let $R^{\mathrm{aff}}$ be the affine row constructed above with full
kernels, $\mu_f$ its completed mean, and
$R_{\mathrm{stat}}=D_f\mu_f+F_f$.  The negative-log vertices are
\[
 \frac b2\bigl(|R_i|^2+|D_iv|^2-|R_i^{\mathrm{aff}}|^2\bigr),
 \quad bU_E,\quad\frac b2|R_{\mathrm{stat},i}|^2,
 \quad\hbox{the old canonical slots},\quad
 -\log\frac{J(u_x)}{J(0)},
\]
with the normalized scalar determinant understood.  This is
\Rref{rg:canonical-vertices} with the physical and added scalar
squares paired before their affine subtraction.  Substitute
$\xi=\mu_f+gZ$.  Assign order $g$ to a relative output tangent and to
$gZ$, and retain total order at most four in the negative logarithm
of the Gaussian expectation, understood here as its formal connected
series at the constant term one.  Equivalently, for each ordered multiset
of at most four vertices, use its connected Wick expectation with
coefficient $(-1)^{v+1}/v!$; internal as well as intervertex pairs are
included.

An anchor is removed by a simultaneous orthogonal transformation.
A smooth local choice of that transformation exists on each fixed
small chart, for example the rotation from its fixed center.  The
resulting tensors are equivariant and do not depend on that local
choice.  At first tangent, the sum of the physical and added scalar
squares agrees with the affine square.  Harmonic completion gives
zero first tangent for the stationary residual.  Hence the row
correction begins with a cubic containing a fluctuation.  The
possible cubic of $U$ pairs a normal quadratic defect with a tangent
first-order row and is zero.  Connected contractions require the
same number of pairs as in \Rref{rg:canonical-contraction}; the
ambient Gaussian has three colors and at most four vertices contribute
through order four.  Odd output invariants vanish by $O(2)$ symmetry.
The kinetic Taylor slots have even parity and unit affine Hessian.

For precision, write the retained polynomial as
$\sum_{r,d}g^{-2d}F_{r,d}(y)$, where $F_{r,d}$ is homogeneous of field
degree $r+2d$ and $r\le4$.  Scalar terms are extracted per unit
volume.  Let $S_{K',o}^{[n]}$ denote the homogeneous degree-$n$ part
of the unmasked output kinetic density, assigned to its anchor by
splitting edge and row terms between their endpoints.  The triangular
prescription is
\begin{align*}
 P'_4&=F_{2,1},&
 \alpha&=\text{affine Hessian of }F_{2,0},&
 I'_2&=\operatorname{Comp}\bigl(F_{2,0}-\alpha S_{K'}^{[2]}\bigr),\\
 P'_5&=0,& I'_3&=0,\\
 P'_6&=F_{4,1},&
 I'_4&=F_{4,0}-\alpha P'_4-\alpha S_{K'}^{[4]},&
 \kappa&=\text{affine Hessian of }F_{4,-1},\\
 J'_2&=\operatorname{Comp}\bigl(F_{4,-1}-\kappa S_{K'}^{[2]}\bigr).
\end{align*}
Here the affine Hessian is the bulk sum per anchor, and
$\operatorname{Comp}$ is the quadratic packet compensation of
Section~\ref{sec:setup}.  Once the total affine coefficient is
removed, the individual compensation coefficients sum to zero, so
this operation preserves the polynomial.  These are the formulas
\Rref{rg:canonical-map} with the odd slots set to zero.  They define a map
$\mathbf P'=\mathcal C_h(\mathbf P)$ that depends only on the free
history and the incoming canonical coefficients.  After the displayed
powers of $g$ have been sorted, neither this map nor its kicks depend
on the precise coupling, cutoffs, regular error, covering gas, or
period.  In the exact output representation we assign the canonical
slots to this finite coefficient operation; its difference from the
actual integral is retained in the error and gas, with the further
affine correction included in $\Delta$.  Thus a run initialized with
$\mathbf P=0$ has exactly the deterministic canonical orbit of these
maps, while its precise coupling obeys Equation~\eqref{eq:orig-4}.
This is the distinction used in the shooting argument.

Coefficient sums are absolutely convergent in the single fixed
coefficient norm.  The inputs to the summability proof in
\Rref{rg:canonical-contraction} and
\Rref{supp:canonical-fixed-norm} are exponential free localization
with its fixed moments, finite-degree tensor contractions and old
exponential path norms.  Our chart and projection vertices have
exactly these properties.  In an off-diagonal contraction, an old
path may be rescaled by the step size and joined by free-kernel
paths.  Only a fixed number of paths occurs per vertex at this order;
their polynomial weights are paid by the old spare exponential or
by kernel moments.  Choose the coefficient exponent $\sigma$ small
for this finite set of moments before choosing $L$.  These estimates
give the asserted $C_L$ bounds and their history and coupling
Lipschitz versions.  Folding and affine compensation use this same
norm, not a succession of norms with smaller exponential exponents.

\begin{lemma}[Dimension-two coefficient contraction]
\label{lem:rg-canonical-contraction}
For the diagonal response of each nonzero canonical slot, the
coefficient norm of the output is at most $CL^{-2}$ times the input
norm, where $C$ is independent of $L$ and the history before the
block size is chosen.  The corresponding quadratic transfer has
zero affine Hessian.  The same conclusion holds for an inclined
step in its rotated physical coordinates.
\end{lemma}

\begin{proof}
For degree $n\ge4$, \Rref{supp:one-norm-contraction} gives
$CL^{2-n}$ and hence the stated bound.  Only the compensated quadratic
slots require an improvement.

At an old anchor $o$, expand the conditional-mean row in a displacement
$r$ through spatial degree two:
\[
 P(o+r,\cdot)-P(o,\cdot)
       =A_o(r)+B_o(r,r)+E_o(r).
\]
Here $A_o$ is homogeneous linear and $B_o$ homogeneous quadratic in
$r$.  In the exponentially weighted row norm of
\Rref{supp:row-second}, their coefficients have bounds $C/L,C/L^2$,
and
\[
 \norm{E_o(r)}\le
 C(1+|r|)^3L^{-3}e^{C\sigma(1+|r|/L)}.
\]
One obtains this formula by a forward Newton polynomial of degree
two, then telescoping its remainder along shortest paths.  The
linear coefficient includes the second-difference adjustment from
rewriting that Newton polynomial in homogeneous powers.  Telescoping
the second and then first differences uses only the third fine
differences supplied by the free estimate.  Thus no continuous
extension of the lattice kernel has been assumed.

Insert the expansion into a compensated inversion-invariant quadratic
packet.  The two-linear-row terms cancel as coefficient arrays by
its zero affine Hessian.  For quarter turns, this uses the symmetric
quadratic polarization fixed above.  The linear--quadratic terms are
odd under $r,s\mapsto-r,-s$ and cancel by the packet's spatial
inversion.  Every remaining term has bound $CL^{-4}$ times a fixed
polynomial in the path length, with its rescaled exponential weight.
The spare old exponential pays that polynomial exactly as in
\Rref{supp:quadratic-before-anchor}.  Output compensation costs a
fixed multiplier.  Summing the $L^2$ old anchors over one new anchor
leaves $CL^{-2}$.

The constant and affine identities for $P$ show that an affine
output tangent transfers to an affine input tangent of slope divided
by the step size.  Each old compensated packet is zero on that
field.  The aggregate transferred affine Hessian is therefore exactly
zero.  Inclined coordinates merely rotate and rescale this argument;
the same row estimates hold in their physical units.
\end{proof}

The canonical calculation has now identified the coefficients and
controlled their own variation.  The remaining obstruction to
Theorem~\ref{thm:rg-step} is an old regular error appearing once,
without another small factor.  Its contraction is the subject of the
next subsection.

\subsection{Contraction of the regular error}
\label{subsec:rg-error}

In the empty primary pattern, the isolated response of an old error
is the operator
\[
 \mathcal T f_X(V)=\E_{C_s}\,[\chi_X(q)f_X(q)],\qquad
 q_x=\operatorname{Exp}_{T_x}
                \bigl(\operatorname{proj}_{T_x^\perp}(\mu_x+gZ_x)\bigr).
\]
Here $\chi_X$ is the prepared smooth graph cutoff, and the output
anchor is the projection of the old anchor.  This is
\Rref{rg:error-transfer} for the intrinsic spin prescription.  For
load $s_X\le L^{1/4}$ its output graph is enlarged to a complete
available good-output ball and all required stencils; the exact
covering correction is retained.  After projection the radius is a
fixed multiple of $M'$.  The entire orbit packet is kept together.

\begin{lemma}[Contraction after removal of the affine Hessian]
\label{lem:rg-error-contraction}
For every fixed required output support exponent chosen before $L$,
\begin{equation}
 \norm{\mathcal Tf}_{k,j-1,B}
   \le\bigl(C/L^2+o_H(1)\bigr)\norm{f}_{k,j,A}.
 \label{eq:orig-9}
\end{equation}
The little-oh may depend on $L$ and is uniform for $H_j\ge H$.
The same bound holds for an input difference through order $k-1$.
A map change acting on an unchanged input satisfies the original
map-difference estimate of \Rref{rg:error-contraction}, with its
polynomial losses and one extra derivative on that input.
\end{lemma}

\begin{proof}
Use an even smooth Gaussian guard on all coordinate reads,
$\max_x|Z_x|\le c_Lp$, with a transition strictly inside the required
plateaus.  The normal component is included.  The prediction lemma,
with output graph chords allowed through $4t'$, puts the unperturbed
fine chords strictly inside the graph cutoff plateau for $L$ large.
Choose $c_L$ small enough that the same holds with the guarded
fluctuation, and on the Taylor segments below.  The guard complement
has an exponentially small cost times fixed polynomials in $M$ and
the load, by the Gaussian union tail and the joint derivative
estimate.  These statements apply on arbitrary graphs, not only
within a single chart.

First suppose $s_X\le L^{1/4}$.  For output derivatives of order at
most three, omit the guarded fluctuation by Taylor expansion twice
in it.  Its first term integrates to zero because the guard and the
Gaussian law are even.  The argument of
\Rref{rg:fluctuation-omission} then gives
\[
 C_LM^D(1+s_X)^Dp^{-2}[f_X]_{k,j}.
\]
It requires at most five old derivatives.  The same argument applies
to the tangent projection in the present spin formula.  Choosing
$P_0$ after the fixed $M$ powers makes this a delayed small factor.

Remove the old anchor by a simultaneous orthogonal transformation
and use geodesic relative coordinates for the unperturbed spins.
The leading first tangent is the conditional-mean prediction.  Expand
it, as above, in $x-o$ into a homogeneous linear spatial polynomial,
a homogeneous quadratic spatial polynomial, and a remainder.  In the
old normalized direction norm, with its low output jets, their bounds
are respectively
\begin{equation}
 C/L,\qquad C/L^2,\qquad C/L^3+o_H(1).
 \label{eq:orig-10}
\end{equation}
Nonlinear chart terms and window truncation belong to the final
remainder.

Here is the weighted-norm verification of this assertion.  If
$|x-o|\le L$, coarse exponential row moments sum both the output
chord variation and the output test-direction weights.  Fine
differences through order three give the stated powers of $L$ with
uniform leading constants.  If $|x-o|>L$, an output direction weight
based at $[x]$ is bounded in the row sum by
$C(1+|x-o|/L)^8$.  Dividing by the old weight
$(1+|x-o|)^8$ gives $O(L^{-8})$; the polynomial predictions based at
$o$ obey the remainder estimate with this same division.  Chord
coordinates have at most linear growth from the anchor in the ball,
so their values satisfy the same argument.  Nonlinear terms in each
fixed jet have an extra $t$ with at worst $C_LM^D$ cost.  Paths outside
the supplied stencils have exponentially small weighted sums.
The ball lies in a common chart because its output radius is $O(M)$
and $M^Dt\to0$.

For completeness, the chain rules can be interpreted entirely in the
norms already defined.  First tangent directions at a spin are tested
in great-circle coordinates, with rotation axes perpendicular to that
spin.  Their symmetric jets are covered by the rotation-derivative
norm.  To differentiate a different chart or a composition, use
normal coordinates at the point of evaluation and the usual chain
rule.  Chart changes on these fixed small charts have bounded fixed
derivatives.  Products of direction weights can occur in one input
slot, including for an unrestricted output rotation axis.  After
one slot's weight is divided out, each additional direction factor
is bounded by a power of
\[
 t\bigl(1+C\mathfrak m_j(1+s_X)\bigr)^8.
\]
On bounded loads this tends to zero with $H$.  At arbitrary load it
has only a fixed polynomial load cost, and each additional mesoscopic
power accompanies an extra $t$.  Also $t'/t$ is uniformly bounded.
Thus the higher map jets asserted to be $o_H(1)$ have that meaning
in the actual graph norm; no uniform bound on all test directions
on arbitrarily large graphs has been assumed.

Expand the invariant error at alignment.  Its constant and first
terms vanish, and its affine Hessian is zero.  Tangent inversion in
the anchor stabilizer makes its cubic term zero.  In the quadratic
term, the two linear spatial predictions vanish by affine
normalization, and the linear--quadratic predictions cancel under
position inversion of the packet.  Equation~\eqref{eq:orig-10} then
gives $CL^{-4}+o_H(1)$ for every surviving quadratic term.  The
quartic Taylor remainder has the same gain.  For the derivatives
through order three, Taylor-expand the required derivative with the
corresponding reduction of order; fourth and, where needed, higher
available old derivatives suffice.  The Taylor segments stay in the
old graph domain: anchor alignment is an isometry and small-log
interpolation distorts chords by a bounded factor tending to one.
The guard leaves a strict margin to the domain boundary.

For output orders four through $k$, retain the fluctuation and apply
the chain rule directly.  Each first derivative of the relative map
costs $C/L+o_H(1)$; differentiating a moving tangent projection of the
fluctuation has an extra $t$.  Every higher map derivative is
$o_H(1)$ in the same norm.  A chain-rule term either has at least four
first-map derivatives or has a higher derivative.  It therefore
again gains $CL^{-4}+o_H(1)$, without requiring an old derivative
above the requested order.  Summing $L^2$ old anchors and the fixed
output support price proves Equation~\eqref{eq:orig-9} on these
short labels.

If $s_X>L^{1/4}$, use the direct composition estimate at the end of
\Rref{rg:error-contraction}.  On the guard, its absolute first jet is
the bounded row of $P_h$, with no bare leading $M$ loss.  Higher jets,
including moving tangent projections, have extra $t$ and fixed
polynomials in $M,1+s_X$.  This statement is local at each halo and
does not require a chart for the whole graph.  Direction-anchor
translations have polynomial load cost; for lifted torus paths the
physical distance obeys the same scaled upper bound.  The spare old
support exponential gives $e^{-cL^{1/4}}$, paying the anchor sum and,
if desired, $L^{-4}$ before it.  The guard-complement estimate has a
bare $M$ power only together with exponential accuracy.  This proves
the estimate also for long and winding labels.

An input difference uses the identical proof in order $k-1$.  For a
changed map, the verified composition and window estimates act on
one additional derivative of the unchanged input, at most $k$.
Thus no derivative is lost from one iteration to the next.  All
arguments retain the graph masks; the strengthened orbit masks and
short-load balls are accompanied by their exact covering corrections.
\end{proof}

\subsection{Taylor comparison, normalization and closure}
\label{subsec:rg-closure}

We now compare the exact connected expression with the canonical
polynomial and close the retained class.  The comparison is an
estimate on the actual integrals; it is not an assumption that a
formal coefficient computation describes them.

An ordinary primitive of order $r>0$ has the majorant established in
Proposition~\ref{prop:rg-joint}.  Inflate it by $g^{-.96r}$ in the
tree condition \Rref{rg:tree-condition}.  The remaining $g^{.04r}$
beats every fixed logarithmic loss, including the $M^2$ site-hit cost.
Thus connected terms of total order at least five retain $g^{4.8}$;
a smaller exponent pays the finite derivative and logarithmic costs.
Old errors count as order four with their stronger actual cap, so
only their isolated response needs Lemma~\ref{lem:rg-error-contraction}.
All other such terms have order at least five.  The resulting
remainder and discrepancy bounds are \Rref{rg:high-order-remainder},
namely $C_Lg^{4.6}$ and $C_Lg^{4.5}$ times the normalized discrepancy.
Exceptional terms, including the new normal-coordinate tail letters,
and canonical paths beyond the mesoscopic cutoff have arbitrary-power
accuracy.

For the terms through order four, group their at most four origins at
one anchor and strengthen their masks to a common good-output ball
containing all reads.  Choose its fixed radius multiplier and then
$L$ large: the old canonical balls shrink by $L$, whereas the fixed
number of new stencils requires only fixed multiples of $M$.
Off-mask terms are retained in the gas.  On this ball and the even
Gaussian guard, prepared factors have their smooth vacuum formulas.
For factors negligible at this accuracy, such as tail letters, that
coincidence is not needed.  Choose angular cutoffs with a large enough
fixed $C_L$ for these plateaus.  The old graph-edge cutoff through
$t/2$ uses the fine-edge prediction, first choosing $L$ large and then
the Gaussian guard constant small.

In these local charts $u,v$ and output coordinates $w_0$ have size
$g\operatorname{poly}(M,p)$.  Every row formula is a composition of
bounded multilinear kernel maps, finitely many uniformly smooth
pointwise functions and their truncated windows; its positional
summations have fixed polynomial weights.  Taylor expansion in
$w_0,gZ$ through field degree six for a vertex with prefactor $g^{-2}$
leaves at least seven small fields.  It therefore costs
$C_Lg^5\operatorname{poly}(M,p)$.  Vertices with other prefactors are
expanded to the corresponding degree.  A normalized output derivative
supplies its factor $t$ and spatial weight.  Taylor expansion with
its degree reduced by the number of derivatives, or a direct smooth
bound once those factors are already supplied, retains the same
$g^5$ estimate through every required order.

The quadratic and harmonic cancellations were established above.
Every expanded exponent starts at its indicated positive order on
the guard, so exponentials minus one and Jacobian logarithms have the
same Taylor control.  Truncated first-jet errors retain their
exponential factor.  The same proof with one marked discrepancy uses
the kernel differences and Proposition~\ref{prop:rg-joint}.  Removing
the guard, restoring full Gaussian polynomial moments, and replacing
short kernels and covariance by full ones have arbitrary-power
accuracy with every fixed positional moment.  Any connection created
only by that replacement contains a long kernel path and has the same
price.

On a fixed origin multiset, the finite product and Mayer-log
identities are identities of these finite polynomials.  They give
exactly the connected Wick prescription above, including repeated
origins and the Jacobian logarithm, as in the proof following
\Rref{rg:high-order-remainder}.  This proves the exact Taylor comparison
and yields \Rref{rg:raw-error} with
$q_1=C/L^2+o_H(1)$.

It remains to return the raw error's constant and affine-Hessian
components to the bulk scalar and kinetic coefficient.  Put
$b^-=b+\alpha+\kappa/b$, and let $R_{\rm raw}$ denote the raw
regular norm after the canonical prefactors have been expressed using
$b^-$.  The preceding estimates give
\[
 R_{\rm raw}\le q_1\delta_j+C_Lg^{4.6},
 \qquad q_1=C/L^2+o_H(1).
\]
For a nonwinding raw label $f_i$ at $o$, fix $n\in S^2$ and set
$v_i=f_i((n)_x)$, its value on the constant spin configuration.
This value is independent of $n$ by $O(3)$ invariance.  Let $b_i$ be
its unit-affine Hessian at that configuration.  The graph norm gives
\[
 |v_i|\le[f_i]_{k,j-1},\qquad
 |b_i|\le C(t')^{-2}[f_i]_{k,j-1}.
\]
There is no support-size loss: a unit affine direction at displacement
$r$ is bounded by the permitted polynomial direction weight.  With
the exact off-mask correction retained, subtract $v_i+b_iS_o$ as in
\Rref{rg:normalization-identity}.  The normalized function has zero
value and affine Hessian and costs only a fixed multiplier in norm.
Summing the bulk coefficients gives
\[
 \Delta=\sum_i b_i,
 \qquad |\Delta|\le C(t')^{-2}R_{\rm raw},
 \qquad b'=b^-+\Delta.
\]
The corresponding kinetic correction is
$-\Delta(\mathcal E'-\sum_oS_o)$.  Allocate $S_o$ to the kinetic
paths in proportion to their affine Hessians, whose sum is one.
Each resulting path expression has zero affine Hessian on its mask,
and \Rref{rg:kinetic-reset} gives total regular norm
$C(t')^2|\Delta|\le CR_{\rm raw}$.  The constant is fixed before
$L$, using uniform free-row moments and the fixed derivative order.
Active $\ell$ rows obey their masked bounds; inactive rows retain
their specified convention.  Mask complements are covering
corrections with the same site-hit estimates.

The last coupling change does not alter the assigned canonical
coefficient lists.  Indeed, retaining their masks, the negative-log
correction for their changed displayed prefactors is exactly
\[
 -\Delta P'_4-\Delta P'_6
       +\bigl((b^-)^{-1}-(b')^{-1}\bigr)J'_2.
\]
Here each polynomial denotes its summed masked list; $I'_2,I'_4$
have no coupling prefactor.  Conversion of the fixed coefficient norm
to $\|\cdot\|_{\rm graph}:=\|\cdot\|_{k,j-1,A}$ on these
regular lists, using its spare path exponential, gives
\[
 \|P'_4\|_{\rm graph}\le C_L(t')^4\operatorname{poly}(M'),
 \quad
 \|P'_6\|_{\rm graph}\le C_L(t')^6\operatorname{poly}(M'),
 \quad
 \|J'_2\|_{\rm graph}\le C_L(t')^2\operatorname{poly}(M').
\]
These bounds hold through order $k$ on the respective graph domains;
path-length and direction-weight powers are summed with the
coefficient exponential.  Since $b^-,b'$ are comparable to $H_j$,
the combined correction has norm at most
\[
 C_LR_{\rm raw}\operatorname{poly}(M')
           \bigl((t')^2+(t')^4+H_j^{-2}\bigr)
       =o_H(1)R_{\rm raw}.
\]
It already has zero value and affine Hessian: $P'_4,P'_6$ have degree
at least four at alignment, and $J'_2$ is compensated.  Thus it can
be retained directly as a regular error on the same masks, without a
further coupling extraction or a change of canonical coordinates.
The marked versions of these products give the same discrepancy
estimates, using $|b_1^{-1}-b_2^{-1}|\le C|b_1-b_2|H_j^{-2}$
for comparable couplings.  This proves quantitatively that the
exact coefficient lists follow $\mathcal C_h$ while the remaining
feedback enters $\Delta$ and the renewed errors.

Consequently \Rref{rg:closed-error} holds with the improved $q_1$,
and affine extraction gives Equation~\eqref{eq:orig-4}.

We give the period argument explicitly because asymmetrical tori are
needed later.  All orbit regroupings use fixed graph enlargements.
A short complete calculation is the corresponding bulk calculation
with its finitely many lifted spin positions refolded.  Its symmetry
is the symmetry of those labels, not a symmetry assertion about the
folded spin configuration.  If a fixed enlargement crosses the
shortest-period threshold, declare the result winding before the
orbit manipulation.  The spare support exponent pays its winding
record, and the lower load retained in graph projection preserves
any winding price already present.  Fold the bulk coefficients and
subtractions with their off-mask corrections.  Missing bulk constants
or affine corrections on a given period are subtracted with exactly
the winding charges of their long derivations.  Actual winding terms
need no bulk constant or Hessian normalization.  Thus the scalar and
coupling increments are independent of period shape, while all
period disagreements retain their stated price.  This is the
mechanism of \Rref{rg:normalization}, with shortest-period length
replacing a square side.

Finally choose the constants.  All support enlargements, derivative
orders and polynomial exponents above are fixed before the scale
choices.  Take $L$ sufficiently large that the fixed multipliers of
$L^{-2}$, including error reset and support dilation, leave a strict
margin below $L^{-2+\upsilon}$.  Choose successive canonical caps and
triangular comparison weights using
Lemma~\ref{lem:rg-canonical-contraction}.  Apply the raw-error
discrepancy and the stronger-than-cap gas discrepancy with their
relative normalizations; error and gas comparison weights may be
reduced if needed.  The input coupling parameter in the cited
comparison estimates is $O_L(|\lambda|/H_j)$.  In the normalized
regular error it incurs only powers of $\log H_j$; affine extraction
multiplies the raw cap by $(t')^{-2}$, a bounded factor after the cap
is inserted.  The high-order and gas discrepancy reserves pay the
remaining powers.  These statements give
Equation~\eqref{eq:orig-5}.

As in \Rref{rg:finite-choices}, choose $P_0$ next, larger than all the
fixed logarithmic requirements, and then $H$.  Any strictly positive
power of $g$ in Equations~\eqref{eq:orig-7} and \eqref{eq:orig-8},
and in the other estimates above, beats the now fixed logarithmic
powers at this last choice.  The little-oh terms are uniform for
$H_j\ge H$.  This completes the proof of
Theorem~\ref{thm:rg-step}.
